% LaTeX resume using res.cls
\documentclass[overlapped,9pt]{res}
\usepackage{url}
%\usepackage{helvetica} % uses helvetica postscript font (download helvetica.sty)
%\usepackage{newcent}   % uses new century schoolbook postscript font
\setlength{\textwidth}{6.0in} % set width of text portion

\begin{document}
%\def\los{{\sl $\ell$OS}\xspace}
\newcommand{\los}{{\sl $\ell$OS}}

% Center the name over the entire width of resume:
 \moveleft.5\hoffset\centerline{\LARGE\bf Evangelos (Vaggelis) Atlidakis}
 \vspace{5pt}
% Draw a horizontal line the whole width of resume:
 %\moveleft\hoffset\vbox{\hrule width\resumewidth height 1pt}\smallskip
% address begins here
% Again, the address lines must be centered over entire width of resume:
 \moveleft.5\hoffset\centerline{400 West 119th Street}
 \moveleft.5\hoffset\centerline{New York, NY 10027}
 \moveleft.5\hoffset\centerline{vatlidak@cs.columbia.edu}
 % \moveleft.5\hoffset\centerline{(205) 566-1765}
 \moveleft.5\hoffset\centerline{\url{http://www.cs.columbia.edu/~vatlidak/}}

\frenchspacing
\begin{resume}


\section{RESEARCH INTERESTS}
  Operating Systems, Distributed Systems, Machine Learning Systems

\section{EDUCATION}
  {\bf Columbia University}, New York, NY \hfill Oct. 2013 - Present\\
  Ph.D. in Computer Science \\
  GPA: 3.9/4 (current)\\
  Advisors: Roxana Geambasu, Jason Nieh \\
  %{\bf Army Service}, Athens, Greece \hfill Aug. 2011 - Aug. 2012 \\
  %Trained with the ``Land Minefield Clearance Battalion''\\
  {\bf University of Athens}, Athens, Greece \hfill Sep. 2005 - Aug. 2011 \\
  Bachelor of Informatics and Telecommunications \\
  Advisors: Alex Delis, Mema Roussopoulos\\
  GPA: 9/10\\
  {\bf Helsinki University of Technology}, Helsinki, Finland \hfill Jan. 2009 - Jun. 2009\\
  Exchange Student, Department of Computer Science \\
  GPA: 10/10

\section{EXPERIENCE}
  % {\sl Graduate Research Assistant} \hfill Oct. 2013 - Present \\
  % Columbia University, New York, NY\\
  % Conducting research in Operating Systems and  Privacy \& Transparency
  {\sl Graduate Teaching Assistant} \hfill Spring 2014 \\
  Columbia University, New York, NY \\
  COMS E6998: Advanced Distributed Systems (Instructor: Roxana Geambasu)


  {\sl Graduate Teaching Assistant} \hfill Fall 2013\\
  Columbia University, New York, NY \\
  COMS W4995: Distributed Systems (Instructor: Roxana Geambasu)

  {\sl Research Associate} \hfill Sep. 2012, May - Sep. 2013 \\
  CERN, Geneva, Switzerland\\
  % Project Associate in CERN, IT Department.\\
  Worked on the Agile Infrastructure project; a IaaS CERN-private cloud\\
  Supervisor: Ignacio Reguero

  {\sl Undergraduate Researcher} \hfill Feb. 2010 - May 2011 \\
  University of Athens, Athens, Greece\\
  Conducted research on distributed systems\\
  Advisors: Alex Delis and Mema Roussopoulos

\section{RESEARCH:}
My research interests are in computer systems, including operating systems
(OSes), distributed systems, and machine learning (ML) systems. Current, I focus
on addressing challenges inherent in the development of modern mobile, cloud,
and ML applications. These applications are fundamentally different from
traditional desktop applications and require renewed support from operating
systems as well as new development environments. I have been running measurement
studies to identify missing or mismatched abstractions for modern applications,
and invent new programming abstractions and systems to support the needs of
modern applications. Following are some example projects in this space.

\textbf{Measurement of POSIX Abstractions in Modern OSes:}\\
In my latest work [1, 2], I demonstrated that the abstractions offered by
traditional OS standards, such as the POSIX API, are insufficient to support
modern applications, such as those running on Android, iOS, and OSX. Modern
applications rely upon very different abstractions, which are currently supplied
by user-space libraries implemented atop POSIX. This layering
causes mismatches, inefficiencies, and even security risks.  For example, I
found that new abstractions typically rely on POSIX extension APIs (i.e., {\tt
ioctl}) to implement their functionality, suggesting that POSIX lacks
appropriate abstractions for modern
workloads. Extension APIs are problematic, because their invocations cannot be
mediated by the OS, putting pressure on user-space libraries and kernel device
drivers to implement correct and coherent protections of these invocations.
I am now studying specific implications and identifying the new OS abstractions
needed to more securely support modern applications.

\textbf{Testing Tools for Data-Driven Applications:}\\
A key aspect characterizing modern applications is their increased reliance on
data and data-driven decision making.  While often beneficial, this practice can
have subtle detrimental consequences, such as discriminatory or racially
offensive effects. I argue that such effects are bugs that should be tested for
and debugged in a manner similar to functionality, reliability, and performance
bugs. To this end, I developed {\em FairTest}, a testing toolkit for data-driven
applications that identifies unwarranted association between an application’s
outputs and user subpopulations, including sensitive groups (e.g., defined by
race or gender). Our paper [3], accepted in the second European Symposium on
Security and Privacy, has attracted attention from several investigative
journalists, who are considering using it to study unwarranted associations in
several applications.

\textbf{Security of Machine Learning Systems under Adversarial Settings:}\\
State of the art machine learning systems, such as those using deep neural
networks for classification and regression, have recently achieved
unprecedented success in a variety of tasks ranging from image and speech
recognition to data compression.
Recent work however has made an intriguing discovery: these systems are
extremely susceptible to adversarial attacks.
That is, an adversary can modify correctly classified samples by adding an
infinitestament amount of carefully crafed noise that will force the model
to misclassify (with high confidence) samples only marginaly different from the
original ones.
In lieu of
these vulnerabilities, I am currently working on defences to harden machine
learning systems and increase their robustness against adversarial attacks.
My long-term goal is to contribute solutions that will shed light
into arcane ML models and will arm developers with powerful debugging primitives
to help identify and address subtle security, performance, and reliability
challenges.

\textbf{Previous Projects:}\\
Before joining Columbia, I worked on several projects related to distributed,
peer-to-peer systems, such as BitTorrent. In one project, I enhanced
BitTorrent with an optimistic unchocking policy that improves the quality of
inter-connections amongst peers and increases the number of peers that are both
directly connected and interested in cooperation. My evaluation showed that the
approach significantly outperforms BitTorrent's original unchocking policy by
(1) increasing the number of directly cooperating peers, (2) easing the load on
seeders by having more peers act as data intermediaries, and (3) shortening the
bootstrapping period for fresh peers [3, 4].


\section{PUBLICATIONS}
  [1] V. Atlidakis, J. Andrus, R. Geambasu, D. Mitropoulos, and J. Nieh.
  ''POSIX has become outdated.'' USENIX ;login: Magazine, 41(3), Fall 2016.

  [2]  V. Atlidakis, J. Andrus, R. Geambasu, D. Mitropoulos, and J. Nieh.
  ``POSIX Abstractions in Modern Operating Systems: The New, the old, and
  the Missing.'' In Proceedings of the eleventh European Conference on Computer
  Systems (EuroSys '16), London, UK, April 2016.

  [3] F. Tramer, V. Atlidakis, R. Geambasu, D. Hsu, J-P Hubaux, M. Humbert,
  A. Juels, and H. Lin. ``Discovering Unwarranted Associations in Data-Driven
  Applications with the FairTest Testing Toolkit.'' Accepted in the second
  European Symposium on Security and Privacy (Euro S\&P '17), Paris, France,
  April 2017.

  [4] V. Atlidakis, M. Roussopoulos, and A. Delis.
  ``EnhancedBit: Unleashing the Potential of the Unchoking Policy in the
  BitTorrent Protocol.'' Journal of Parallel and Distributed Computing (JPDC),
  2014.

  [5] V. Atlidakis, M. Roussopoulos, and A. Delis.
  ``Changing the Unchoking Policy for an Enhanced Bittorrent.'' International
  European Conference on Parallel and Distributed Computing (Euro-Par), 2012.

\section{HONORS}
  \vspace{5pt}
  Entered department of Informatics and Telecommunications with highest score, Sep. 2005.\\

  \vspace{-20pt}
  GPA in top 3\% of department of Informatics and Telecommunications, Dec. 2011.

\section{TECHNICAL \\ SKILLS}

  \vspace{5pt}
  {\sl Operating Systems:} Ubuntu, Debian, RHEL, CentOS \\

  \vspace{-20pt}
  {\sl Languages:} C/C++, Python, Java, Shell Scripting\\

  \vspace{-20pt}
  {\sl Systems/Frameworks:} Tensorflow, Theano, IBM CPLEX, Linux Kernel,
    Android\\

  \vspace{-20pt}
  {\sl Misc:} Matlab, IDA, Apache, Nagios, Puppet, Foreman, OpenStack


\section{COURSES} Operating Systems, Advanced Operating Systems,
    Programming Languages, Machine Learning, Natural Language Processing,
    Algorithms, Cryptography.


\end{resume}
\end{document}
